\chapter{Systèmes d'exploitation, concurrence et réseaux}\label{ch:iv3:operating-systems-concurrency-and-networks}

\begin{ivtrains}
Ce qu'est réellement un processus, un thread et une page, comment un verrou atteint le noyau, et ce qui arrive à un paquet sur son chemin jusqu'au code de l'application. Tout ce chapitre se rapporte à Linux sur x86-64. Les programmes réseau tournent sur la boucle locale, en Python ou en C++20. Il faut savoir quel comportement relève du protocole, lequel relève de Linux et lequel relève de la machine que l'on a sous les yeux.
\end{ivtrains}

\begin{question}[$\star$ \iqroles{QD, Dev, LLE} \iqfirm{toutes} \iqlevel{junior}]\label{iq:iv3:operating-systems-concurrency-and-networks:1}
Sous Linux, quelles deux lignes le \cref{lst:iv3:os:forkthread} affiche-t-il~? Que partage un thread avec les autres threads de son processus, et que possède chacun en propre~? Qu'est-ce qui est partagé après un \texttt{fork}~?
\omcode{code/interviews-3/05-operating-systems-concurrency-and-networks/python/fork_vs_thread.py}{1}{22}{Un processus enfant et un thread.\protect\hfil\protect\break}\label{lst:iv3:os:forkthread}
\end{question}

\begin{question}[$\star$ \iqroles{QD, Dev, LLE} \iqfirm{toutes} \iqlevel{junior}]\label{iq:iv3:operating-systems-concurrency-and-networks:2}
Qu'affiche le \cref{lst:iv3:os:tls}~? Où vit la copie de \texttt{mine} de chaque thread, comparée à \texttt{shared}~? Pourquoi le dernier nombre sort-il tel qu'il sort~? Citez un usage du stockage local au thread pour un compteur lu rarement.
\omcode{code/interviews-3/05-operating-systems-concurrency-and-networks/cpp/tls.cpp}{1}{22}{Une variable par thread.\protect\hfil\protect\break}\label{lst:iv3:os:tls}
\end{question}

\begin{question}[$\star$ \iqroles{QD, Dev, LLE} \iqfirm{toutes} \iqlevel{junior}]\label{iq:iv3:operating-systems-concurrency-and-networks:3}
Quelles trois lignes le \cref{lst:iv3:os:zombie} affiche-t-il~? Qu'est-ce qu'un zombie, ce qu'il occupe encore, et comment s'en débarrasser~? Que devient-il si le parent meurt en premier~?
\omcode{code/interviews-3/05-operating-systems-concurrency-and-networks/python/zombie.py}{1}{18}{Un enfant qui a terminé.\protect\hfil\protect\break}\label{lst:iv3:os:zombie}
\end{question}

\begin{question}[$\star$ \iqroles{QD, Dev, LLE} \iqfirm{toutes} \iqlevel{junior}]\label{iq:iv3:operating-systems-concurrency-and-networks:4}
Quelles deux lignes le \cref{lst:iv3:os:boundaries} affiche-t-il~? Combien d'appels \texttt{recv} la boucle TCP doit-elle effectuer avant de s'arrêter, et ce nombre est-il fixe~? Énoncez ce que TCP garantit et ce qu'UDP ne garantit pas, ainsi que la seule chose qu'UDP conserve et que TCP ne conserve pas.
\omcode{code/interviews-3/05-operating-systems-concurrency-and-networks/python/boundaries.py}{1}{19}{Deux envois sur chaque protocole.\protect\hfil\protect\break}\label{lst:iv3:os:boundaries}
\end{question}

\begin{question}[$\star$ \iqroles{QD, Dev, LLE} \iqfirm{toutes} \iqlevel{junior}]\label{iq:iv3:operating-systems-concurrency-and-networks:5}
Sur x86-64, quelles trois lignes le \cref{lst:iv3:os:endian} affiche-t-il~? Qu'est-ce que l'ordre des octets du réseau~? Un flux envoie un prix sur 32~bits en little-endian~: lequel des deux formats se lit directement sur cette machine, et que faut-il sur une machine big-endian~?
\omcode{code/interviews-3/05-operating-systems-concurrency-and-networks/python/endian.py}{1}{8}{Ordre des octets.\protect\hfil\protect\break}\label{lst:iv3:os:endian}
\end{question}

\begin{question}[$\star$ \iqroles{QD, Dev, LLE} \iqfirm{toutes} \iqlevel{junior}]\label{iq:iv3:operating-systems-concurrency-and-networks:6}
Quelles trois lignes le \cref{lst:iv3:os:fivetuple} affiche-t-il~? Deux clients sont connectés à un même port serveur~: comment le noyau sait-il à quelle connexion appartient un segment~? Qu'est-ce qui identifie une connexion TCP~?
\omcode{code/interviews-3/05-operating-systems-concurrency-and-networks/python/five_tuple.py}{1}{12}{Deux connexions vers un même port.\protect\hfil\protect\break}\label{lst:iv3:os:fivetuple}
\end{question}

\begin{question}[$\star$ \iqroles{QD, Dev, LLE} \iqfirm{toutes} \iqlevel{junior}]\label{iq:iv3:operating-systems-concurrency-and-networks:7}
Avec g++ 11.4, quelles deux lignes le \cref{lst:iv3:os:semaphore} affiche-t-il~? Quelle est la différence entre un sémaphore et un mutex, sachant que \texttt{other} a libéré un emplacement qu'il n'avait jamais acquis~? Que serait ce même geste sur un \texttt{std::mutex}~?
\omcode{code/interviews-3/05-operating-systems-concurrency-and-networks/cpp/semaphore.cpp}{1}{13}{Trois emplacements.\protect\hfil\protect\break}\label{lst:iv3:os:semaphore}
\end{question}

\begin{question}[$\star$ \iqroles{QD, Dev, LLE} \iqfirm{toutes} \iqlevel{junior}]\label{iq:iv3:operating-systems-concurrency-and-networks:8}
Quelles trois lignes le \cref{lst:iv3:os:signals} affiche-t-il~? Pourquoi un processus peut-il intercepter \texttt{SIGTERM} mais pas \texttt{SIGKILL}~? Comment faut-il donc arrêter un processus de trading qui détient des ordres ouverts~?
\omcode{code/interviews-3/05-operating-systems-concurrency-and-networks/python/signals.py}{1}{18}{Deux signaux.\protect\hfil\protect\break}\label{lst:iv3:os:signals}
\end{question}

\begin{question}[$\star\star$ \iqroles{QD, Dev, LLE} \iqfirm{toutes} \iqlevel{junior, confirmé}]\label{iq:iv3:operating-systems-concurrency-and-networks:9}
Le \Cref{lst:iv3:os:cow} touche 2000 pages de 4~KiB, puis fait un fork. Combien de défauts de page l'enfant subit-il à peu près en lisant toutes ces pages, puis en les écrivant toutes~? Quelles valeurs chaque processus affiche-t-il ensuite, et comment les deux peuvent-ils afficher la même adresse~? Expliquez la copie sur écriture, et ce qu'elle implique pour un grand processus qui fait un fork afin de prendre un instantané pendant qu'il continue de mettre à jour ses données.
\omcode{code/interviews-3/05-operating-systems-concurrency-and-networks/cpp/cow.cpp}{8}{37}{Lecture et écriture après un fork.\protect\hfil\protect\break}\label{lst:iv3:os:cow}
\end{question}

\begin{question}[$\star\star$ \iqroles{QD, Dev, LLE} \iqfirm{toutes} \iqlevel{junior, confirmé}]\label{iq:iv3:operating-systems-concurrency-and-networks:10}
Le x86-64 avec une pagination à 4 niveaux et des pages de 4~KiB traduit des adresses virtuelles de 48~bits, avec 512 entrées de 8~octets dans chaque table. Comment l'adresse \texttt{0x7f1234567abc} se découpe-t-elle, et quels sont les quatre indices de table et le décalage, en décimal~? Quelle quantité de mémoire les tables de pages occupent-elles pour mettre en correspondance 1~GiB de mémoire virtuelle contiguë et alignée~? Que change-t-il avec des pages de 2~MiB~? Écrivez-le en Python.
\end{question}

\begin{question}[$\star\star$ \iqroles{QD, Dev, LLE} \iqfirm{toutes} \iqlevel{confirmé}]\label{iq:iv3:operating-systems-concurrency-and-networks:11}
Avec glibc sur Linux x86-64, quelles trois lignes le \cref{lst:iv3:os:mallocmmap} affiche-t-il, sachant que \texttt{where} indique le mappage qui contient une adresse~? Pourquoi un bloc de 1~MiB vit-il ailleurs qu'un bloc de 1~KiB~? Pourquoi la dernière ligne diffère-t-elle du deuxième mot de la première~? Pourquoi un programme sensible à la latence s'en soucierait-il~?
\omcode{code/interviews-3/05-operating-systems-concurrency-and-networks/cpp/malloc_mmap.cpp}{20}{30}{Où malloc obtient sa mémoire. Une fonction d'aide, \texttt{where}, n'est pas montrée.\protect\hfil\protect\break}\label{lst:iv3:os:mallocmmap}
\end{question}

\begin{question}[$\star\star$ \iqroles{QD, Dev, LLE} \iqfirm{toutes} \iqlevel{junior, confirmé}]\label{iq:iv3:operating-systems-concurrency-and-networks:12}
Le \Cref{lst:iv3:os:rssvsz} rapporte la taille virtuelle et la taille résidente, en MiB, à partir de \texttt{/proc/self/status}. Quelles deux lignes affiche-t-il à peu près~? Comment un \texttt{mmap} de 1~GiB peut-il réussir sans utiliser un gigaoctet~? Que se passe-t-il si chaque processus touche à tout ce qu'il a demandé~?
\omcode{code/interviews-3/05-operating-systems-concurrency-and-networks/python/rss_vsz.py}{1}{23}{Demander de la mémoire et l'utiliser.\protect\hfil\protect\break}\label{lst:iv3:os:rssvsz}
\end{question}

\begin{question}[$\star\star$ \iqroles{QD, Dev, LLE} \iqfirm{toutes} \iqlevel{junior}]\label{iq:iv3:operating-systems-concurrency-and-networks:13}
Le \Cref{lst:iv3:os:stackdepth} est compilé avec g++ 11.4, \texttt{-O2}, et exécuté avec \texttt{ulimit -s 8192}, si bien que la limite de pile est de 8~MiB. Chaque appel utilise un peu plus de 1~KiB de pile. Que se passe-t-il avec l'argument 1000, et avec 100\,000~? Et avec 100\,000 sous \texttt{ulimit -s 262144}~? Peut-on intercepter cet échec avec \texttt{try} et \texttt{catch}~? Quelle est la taille de la pile d'un nouveau thread sous la première limite~?
\omcode{code/interviews-3/05-operating-systems-concurrency-and-networks/cpp/stack_depth.cpp}{1}{12}{Descendre dans la pile.\protect\hfil\protect\break}\label{lst:iv3:os:stackdepth}
\end{question}

\begin{question}[$\star\star$ \iqroles{QD, Dev, LLE} \iqfirm{crypto, teneur de marché, hedge fund} \iqlevel{junior, confirmé}]\label{iq:iv3:operating-systems-concurrency-and-networks:14}
Dans le \cref{lst:iv3:os:pingpong}, deux processus épinglés sur un même cœur se renvoient un octet 1000 fois à travers des tubes. Quelle est la plus petite valeur qu'il puisse afficher, et pourquoi~? Qu'est-ce qu'un changement de contexte, et que coûte-t-il au-delà du temps du changement lui-même~? Pourquoi le même ping-pong sur deux cœurs avec attente active évite-t-il cela~?
\omcode{code/interviews-3/05-operating-systems-concurrency-and-networks/python/ping_pong.py}{1}{27}{Compter les changements de contexte.\protect\hfil\protect\break}\label{lst:iv3:os:pingpong}
\end{question}

\begin{question}[$\star\star$ \iqroles{QD, Dev} \iqfirm{toutes} \iqlevel{junior}]\label{iq:iv3:operating-systems-concurrency-and-networks:15}
Trois tâches arrivent à l'instant 0, dans l'ordre A, B, C, avec des besoins en CPU de 8, 4 et 1. Calculez le temps d'achèvement de chacune et la moyenne avec un tourniquet de quantum 2, puis avec le premier arrivé premier servi, puis avec le plus court d'abord. Ne tenez pas compte du coût des changements de contexte. Écrivez le tourniquet en Python. Que devient-il quand le quantum augmente~? Pourquoi ne pas toujours choisir la tâche la plus courte~?
\end{question}

\begin{question}[$\star\star$ \iqroles{QD, Dev} \iqfirm{toutes} \iqlevel{junior, confirmé}]\label{iq:iv3:operating-systems-concurrency-and-networks:16}
On vous indique quel thread détient chaque verrou, et sur quel verrou chaque thread est bloqué (au plus un). En Python, renvoyez une liste de threads formant un interblocage, ou \texttt{None}. Quelles quatre conditions doivent toutes être réunies pour qu'il y ait un interblocage, et laquelle est généralement supprimée en pratique~? Pourquoi cela peut-il se faire sans recherche générale de cycle~?
\end{question}

\begin{question}[$\star\star$ \iqroles{QD, Dev, LLE} \iqfirm{toutes} \iqlevel{junior, confirmé}]\label{iq:iv3:operating-systems-concurrency-and-networks:17}
Sous Linux, quelles deux lignes le \cref{lst:iv3:os:mmapfile} affiche-t-il~? Lequel de ces résultats le manuel promet-il, et lequel n'est-il que ce qui se passe ici~? Quand lirait-on un fichier à travers un mappage plutôt qu'avec \texttt{read}~? Quand l'écriture d'un mappage partagé atteint-elle le disque~?
\omcode{code/interviews-3/05-operating-systems-concurrency-and-networks/python/mmap_file.py}{1}{15}{Un mappage partagé et un mappage privé d'un même fichier.\protect\hfil\protect\break}\label{lst:iv3:os:mmapfile}
\end{question}

\begin{question}[$\star\star$ \iqroles{QD, Dev} \iqfirm{toutes} \iqlevel{junior, confirmé}]\label{iq:iv3:operating-systems-concurrency-and-networks:18}
Un processus sauvegarde son état en écrasant un fichier, et \texttt{write} a renvoyé un succès au moment où le courant a été coupé. Qu'est-ce qui peut se trouver sur le disque~? En Python sous Linux, écrivez \texttt{atomic\_save(path, data)} de sorte qu'un plantage laisse soit l'ancien fichier entier, soit le nouveau fichier entier. Expliquez chaque appel système. Qu'est-ce qui reste non protégé, malgré tout~?
\end{question}

\begin{question}[$\star\star$ \iqroles{QD, Dev, LLE} \iqfirm{crypto, teneur de marché, hedge fund} \iqlevel{junior, confirmé}]\label{iq:iv3:operating-systems-concurrency-and-networks:19}
Le \Cref{lst:iv3:os:clocks} affiche la résolution de \texttt{CLOCK\_MONOTONIC}, et le nombre de fois où elle a reculé sur 200\,000 lectures. Quelles sont les deux lignes sur cette machine~? Vous devez chronométrer une requête et aussi horodater un ordre avec l'heure du jour. Quelle horloge utilisez-vous pour chacun, et pourquoi~? Comment une lecture d'horloge peut-elle ne faire aucun appel système~? Pourquoi la résolution n'est-elle pas la même chose que la précision~?
\omcode{code/interviews-3/05-operating-systems-concurrency-and-networks/python/clocks.py}{1}{10}{Une horloge monotone.\protect\hfil\protect\break}\label{lst:iv3:os:clocks}
\end{question}

\begin{question}[$\star\star$ \iqroles{QD, Dev, LLE} \iqfirm{toutes} \iqlevel{junior}]\label{iq:iv3:operating-systems-concurrency-and-networks:20}
Quelles deux lignes le \cref{lst:iv3:os:acceptlater} affiche-t-il, et dans quel ordre~? Le client termine \texttt{connect} et envoie des données avant que le serveur ait appelé \texttt{accept}. Comment cela est-il possible~? Décrivez la poignée de main, et ce que contrôle l'argument backlog. Que devient un nouveau client une fois la file pleine~?
\omcode{code/interviews-3/05-operating-systems-concurrency-and-networks/python/accept_later.py}{1}{8}{Connexion avant accept.\protect\hfil\protect\break}\label{lst:iv3:os:acceptlater}
\end{question}

\begin{question}[$\star\star$ \iqroles{QD, Dev, LLE} \iqfirm{toutes} \iqlevel{junior, confirmé}]\label{iq:iv3:operating-systems-concurrency-and-networks:21}
Votre serveur redémarre et échoue avec \texttt{Address already in use}, alors qu'aucun processus ne détient le port. Quelles trois lignes le \cref{lst:iv3:os:timewait} affiche-t-il sous Linux~? Quel côté entre dans \texttt{TIME\_WAIT}, et pourquoi cet état existe-t-il~? Pourquoi la deuxième ligne ne dit-elle pas \texttt{bound}, alors que la nouvelle socket a demandé \texttt{SO\_REUSEADDR}~?
\omcode{code/interviews-3/05-operating-systems-concurrency-and-networks/python/time_wait.py}{1}{34}{Redémarrage d'un serveur. L'état 06 est \texttt{TIME\_WAIT}.\protect\hfil\protect\break}\label{lst:iv3:os:timewait}
\end{question}

\begin{question}[$\star\star$ \iqroles{QD, Dev} \iqfirm{toutes} \iqlevel{junior, confirmé}]\label{iq:iv3:operating-systems-concurrency-and-networks:22}
Une liaison de 10~Gb/s a un temps aller-retour de 1~ms, et rien n'est perdu. Quel est le produit bande passante-délai~? Avec une fenêtre de 65\,535 octets, quel débit obtient-on, et comment le facteur d'échelle de fenêtre (window scaling) corrige-t-il cela~? Quelle est la différence entre contrôle de flux et contrôle de congestion~? Énoncez ensuite ce qui arrive au débit lorsque le temps aller-retour double, pour une fenêtre fixe.
\end{question}

\begin{question}[$\star\star$ \iqroles{QD, Dev, LLE} \iqfirm{toutes} \iqlevel{junior, confirmé}]\label{iq:iv3:operating-systems-concurrency-and-networks:23}
Des messages sont envoyés sur TCP avec une longueur de 2~octets en gros-boutiste, suivie du corps. En Python, écrivez un \texttt{Deframer} dont \texttt{feed(data)} prend les octets à mesure qu'ils arrivent de \texttt{recv} et renvoie les corps de messages complets. Rejetez une longueur supérieure à 4096. Pourquoi \texttt{msg = sock.recv(4096)} passe-t-il un test en boucle locale puis échoue-t-il en production~? Qu'en est-il du côté émetteur~?
\end{question}

\begin{question}[$\star\star$ \iqroles{QD, Dev, LLE} \iqfirm{crypto, teneur de marché, hedge fund} \iqlevel{junior}]\label{iq:iv3:operating-systems-concurrency-and-networks:24}
Sous Linux, quelles trois lignes affiche le \cref{lst:iv3:os:udptrunc}~? Que sont devenus les octets \texttt{efgh}~? D'où vient le nombre 65\,507~? Quelle doit être la taille des datagrammes d'un vrai flux, et pourquoi~?
\omcode{code/interviews-3/05-operating-systems-concurrency-and-networks/python/udp_trunc.py}{1}{14}{Datagrammes et tampons.\protect\hfil\protect\break}\label{lst:iv3:os:udptrunc}
\end{question}

\begin{question}[$\star\star$ \iqroles{QD, Dev, LLE} \iqfirm{crypto, teneur de marché, hedge fund} \iqlevel{junior, confirmé}]\label{iq:iv3:operating-systems-concurrency-and-networks:25}
Sous Linux, quelles quatre lignes affiche le \cref{lst:iv3:os:udpdrop}~? Pourquoi l'option renvoie-t-elle une valeur différente de celle qui a été fixée~? L'émetteur n'a signalé aucune erreur, alors qui a perdu les données, et comment le savoir~? Que feriez-vous pour un gestionnaire de flux (feed handler) qui perd des paquets lors des rafales~?
\omcode{code/interviews-3/05-operating-systems-concurrency-and-networks/python/udp_drop.py}{1}{22}{Un récepteur qui ne lit pas.\protect\hfil\protect\break}\label{lst:iv3:os:udpdrop}
\end{question}

\begin{question}[$\star\star$ \iqroles{QD, Dev, LLE} \iqfirm{crypto, teneur de marché, hedge fund} \iqlevel{junior, confirmé}]\label{iq:iv3:operating-systems-concurrency-and-networks:26}
Qu'est-ce qu'un groupe multicast IPv4, et comment un hôte y adhère-t-il~? Exécutez le \cref{lst:iv3:os:mcast} avec \texttt{join}, puis avec n'importe quel autre argument~: qu'affiche chacun sur cette machine Linux~? À quelle adresse Ethernet le groupe 239.1.2.3 correspond-il, et quels autres groupes partagent-ils cette adresse~? Quelles sont les valeurs par défaut de \texttt{IP\_MULTICAST\_TTL} et de \texttt{IP\_MULTICAST\_LOOP}, et pourquoi cela compte-t-il lorsqu'un test reçoit des données sur une machine mais pas entre deux~?
\omcode{code/interviews-3/05-operating-systems-concurrency-and-networks/python/mcast_join.py}{1}{20}{Adhésion à un groupe, sur l'interface de boucle locale.\protect\hfil\protect\break}\label{lst:iv3:os:mcast}
\end{question}

\begin{question}[$\star\star$ \iqroles{QD, Dev, LLE} \iqfirm{toutes} \iqlevel{junior, confirmé}]\label{iq:iv3:operating-systems-concurrency-and-networks:27}
Quelles sont les deux lignes affichées par le \cref{lst:iv3:os:nonblock}, où chacune correspond à ce que \texttt{connect\_ex} a renvoyé, puis à ce que \texttt{SO\_ERROR} contenait~? Pourquoi l'erreur d'une connexion refusée n'est-elle pas en première position~? Pourquoi faut-il vérifier \texttt{SO\_ERROR} lorsque \texttt{select} indique que la socket est inscriptible~?
\omcode{code/interviews-3/05-operating-systems-concurrency-and-networks/python/nonblock.py}{1}{20}{Une connexion qui n'attend pas.\protect\hfil\protect\break}\label{lst:iv3:os:nonblock}
\end{question}

\begin{question}[$\star\star$ \iqroles{QD, Dev, LLE} \iqfirm{toutes} \iqlevel{junior, confirmé}]\label{iq:iv3:operating-systems-concurrency-and-networks:28}
Que se passe-t-il pour le \cref{lst:iv3:os:sigpipe} exécuté sans argument, puis avec un argument~? Dans le \cref{lst:iv3:os:firstwrite}, le correspondant a déjà fermé une connexion TCP~: pourquoi le premier envoi réussit-il, et que fait le second~? Comment un serveur doit-il se protéger contre cela~?
\omcode{code/interviews-3/05-operating-systems-concurrency-and-networks/cpp/sigpipe.cpp}{7}{14}{Écriture vers un correspondant fermé, en C++.\protect\hfil\protect\break}\label{lst:iv3:os:sigpipe}
\omcode{code/interviews-3/05-operating-systems-concurrency-and-networks/python/first_write.py}{1}{15}{Écriture vers un correspondant fermé sur TCP, en Python.\protect\hfil\protect\break}\label{lst:iv3:os:firstwrite}
\end{question}

\begin{question}[$\star\star\star$ \iqroles{QD, Dev, LLE} \iqfirm{toutes} \iqlevel{confirmé}]\label{iq:iv3:operating-systems-concurrency-and-networks:29}
Un processeur, des priorités strictes. Une tâche de faible priorité prend un verrou à l'instant 0 et le conserve pendant 4 ticks de son propre temps CPU. Une tâche de haute priorité arrive à l'instant 1 et a besoin du verrou pendant 1 tick, puis s'exécute 1 tick de plus. Une tâche de priorité moyenne arrive à l'instant 2 et demande 20 ticks, sans verrou. À quel instant la tâche de haute priorité se termine-t-elle~? Comment appelle-t-on ce phénomène~? Que change l'héritage de priorité, et quel est alors l'instant de fin~? Écrivez une simulation tick par tick en Python, et montrez comment le retard dépend de la tâche moyenne.
\end{question}

\begin{question}[$\star\star\star$ \iqroles{QD, Dev, LLE} \iqfirm{toutes} \iqlevel{confirmé}]\label{iq:iv3:operating-systems-concurrency-and-networks:30}
En C++20, à l'aide d'un seul \texttt{std::mutex} et de variables de condition, écrivez un verrou lecteurs-rédacteurs dans lequel un rédacteur en attente bloque les nouveaux lecteurs. Incluez \texttt{try\_lock\_shared}. Pourquoi cette priorité, et quel risque cela comporte-t-il~? Comment tester que les lecteurs ne dépassent pas un rédacteur en attente, sans dépendre du temps~? Pourquoi un verrou lecteurs-rédacteurs pourrait-il n'être pas plus rapide qu'un mutex pour des sections critiques courtes~?
\end{question}

\begin{question}[$\star\star\star$ \iqroles{QD, Dev, LLE} \iqfirm{crypto, teneur de marché, hedge fund} \iqlevel{confirmé}]\label{iq:iv3:operating-systems-concurrency-and-networks:31}
En C++20 sous Linux, écrivez un mutex sur un seul \texttt{std::atomic<int>} et l'appel système \texttt{futex}, de sorte qu'un verrou dont personne d'autre ne veut ne fasse aucun appel système. Pourquoi faut-il trois états, et non deux~? Que se passe-t-il si l'état revient à 1 après le réveil d'un dormeur~? Pourquoi \texttt{FUTEX\_WAIT} doit-il recevoir la valeur attendue en argument~?
\end{question}

\begin{question}[$\star\star\star$ \iqroles{QD, Dev, LLE} \iqfirm{crypto, teneur de marché, hedge fund} \iqlevel{junior, confirmé}]\label{iq:iv3:operating-systems-concurrency-and-networks:32}
Un client effectue deux petites écritures sur une connexion TCP, un en-tête puis un corps, et attend ensuite la réponse. Le temps aller-retour est de 0.1~ms, et le serveur retient un accusé de réception pendant 40~ms au maximum, sauf s'il a des données à envoyer, et il ne répond qu'une fois qu'il a reçu les deux morceaux. L'algorithme de Nagle est activé. Quel est le délai avant l'arrivée de la réponse~? Expliquez le mécanisme. Donnez trois façons de corriger le problème, et dites laquelle est la meilleure. Modélisez le phénomène en Python.
\end{question}

\begin{question}[$\star\star\star$ \iqroles{LLE, QD, Dev} \iqfirm{crypto, teneur de marché, hedge fund} \iqlevel{confirmé}]\label{iq:iv3:operating-systems-concurrency-and-networks:33}
Un paquet de données de marché UDP arrive sur la carte réseau d'un serveur Linux, et votre processus est bloqué dans \texttt{recv}. Décrivez le chemin, étape par étape, jusqu'au retour de \texttt{recv}. À quelles étapes le paquet peut-il attendre, être abandonné ou être copié~? Quelles étapes le contournement du noyau (kernel bypass) supprime-t-il~? Où chercher les pertes à chaque étage~?
\end{question}

\begin{question}[$\star\star\star$ \iqroles{QD, Dev, LLE} \iqfirm{toutes} \iqlevel{confirmé}]\label{iq:iv3:operating-systems-concurrency-and-networks:34}
Quelles sont les deux lignes affichées par le \cref{lst:iv3:os:epoll}~? Expliquez chaque nombre. Pourquoi le mode déclenché par front exige-t-il une socket non bloquante et une lecture jusqu'à \texttt{EAGAIN}~? Pourquoi \texttt{epoll} est-il préférable à \texttt{select} ou \texttt{poll} pour dix mille connexions majoritairement inactives, et dans quel cas ce gain n'existe-t-il pas réellement~?
\omcode{code/interviews-3/05-operating-systems-concurrency-and-networks/python/epoll_modes.py}{1}{21}{Niveau et front.\protect\hfil\protect\break}\label{lst:iv3:os:epoll}
\end{question}

\begin{question}[$\star\star\star$ \iqroles{QD, Dev, LLE} \iqfirm{toutes} \iqlevel{confirmé}]\label{iq:iv3:operating-systems-concurrency-and-networks:35}
Avec g++ 11.4 sous Linux, quelles trois lignes affiche le \cref{lst:iv3:os:forklock}~? Expliquez les deux valeurs de l'enfant. Supposons que l'autre thread ait tenu le verrou à l'intérieur d'une bibliothèque de journalisation, et que l'enfant appelle maintenant cette bibliothèque~: que se passe-t-il~? Que peut faire sans danger un enfant d'un processus multithreadé avant \texttt{exec}~?
\omcode{code/interviews-3/05-operating-systems-concurrency-and-networks/cpp/fork_lock.cpp}{9}{38}{Fork pendant qu'un verrou est tenu.\protect\hfil\protect\break}\label{lst:iv3:os:forklock}
\end{question}

\begin{question}[$\star\star\star$ \iqroles{QD, Dev, LLE} \iqfirm{crypto, teneur de marché, hedge fund} \iqlevel{confirmé}]\label{iq:iv3:operating-systems-concurrency-and-networks:36}
Deux processus partagent un \texttt{pthread\_mutex\_t} en mémoire partagée, et l'un meurt en le détenant. Le \Cref{lst:iv3:os:robust} affiche le résultat du verrouillage par l'autre~: qu'affiche-t-il sans argument, puis avec un argument~? Pourquoi le second cas est-il si dangereux sans limite de temps~? Que doit faire le nouveau propriétaire sur \texttt{EOWNERDEAD}, et pourquoi la bibliothèque ne peut-elle pas le faire à sa place~?
\omcode{code/interviews-3/05-operating-systems-concurrency-and-networks/cpp/robust.cpp}{9}{39}{Un verrou dont le propriétaire meurt.\protect\hfil\protect\break}\label{lst:iv3:os:robust}
\end{question}

\begin{question}[$\star\star\star$ \iqroles{QD, Dev, LLE} \iqfirm{crypto, teneur de marché, hedge fund} \iqlevel{junior, confirmé}]\label{iq:iv3:operating-systems-concurrency-and-networks:37}
Un flux envoie les mêmes paquets, numérotés à partir de 1, sur deux lignes UDP, A et B, de sorte qu'un paquet perdu sur l'une peut arriver sur l'autre. En Python, écrivez un \texttt{Arbiter} dont \texttt{on\_packet(seq)} prend chaque arrivée, de l'une ou l'autre ligne, et renvoie les numéros désormais livrables, dans l'ordre, chacun exactement une fois. Si plus de \texttt{hold} paquets ultérieurs attendent derrière un paquet manquant, déclarez un trou et poursuivez. Pourquoi ne pas attendre indéfiniment~? Comment choisir \texttt{hold}~?
\end{question}

\begin{question}[$\star\star\star\star$ \iqroles{QD, Dev, LLE} \iqfirm{crypto, teneur de marché, hedge fund} \iqlevel{confirmé}]\label{iq:iv3:operating-systems-concurrency-and-networks:38}
En Python, avec \texttt{select.epoll} et un seul thread, écrivez un serveur d'écho pour des messages encadrés par une longueur de 2~octets. Il doit tenir face à des messages coupés à n'importe quel octet, à un client qui écrit plus vite qu'il ne lit, et à un client qui ferme la connexion alors que des réponses sont en attente. Où vont les réponses en attente~? Quand demande-t-on \texttt{EPOLLOUT}, et pourquoi pas systématiquement~? Que fait un client lent aux autres~? Supposez que vous disposez déjà de \texttt{Deframer}, dont \texttt{feed(data)} renvoie une liste de corps complets et lève \texttt{ValueError} pour une longueur supérieure à 4096, et de \texttt{frame(body)}, qui renvoie la longueur et le corps.
\end{question}

\begin{question}[$\star\star\star\star$ \iqroles{QD, Dev} \iqfirm{crypto, teneur de marché, hedge fund} \iqlevel{confirmé}]\label{iq:iv3:operating-systems-concurrency-and-networks:39}
Construisez une livraison fiable et ordonnée sur une liaison non fiable en Python, sous forme de modèle manche par manche~: répétition sélective, avec une fenêtre de $w$ paquets. La liaison peut perdre, dupliquer ou réordonner des paquets, et peut perdre des accusés de réception. Pourquoi le récepteur doit-il acquitter un doublon~? Pourquoi une fenêtre est-elle nécessaire~? Démontrez qu'elle livre dans l'ordre et exactement une fois. Comment la tester, y compris avec un taux de perte de 1~? Comparez avec TCP.
\end{question}

\begin{question}[$\star\star\star\star$ \iqroles{QD, Dev, LLE} \iqfirm{crypto, teneur de marché, hedge fund} \iqlevel{confirmé}]\label{iq:iv3:operating-systems-concurrency-and-networks:40}
Un processus producteur et un processus consommateur partagent un tampon circulaire (ring buffer) en mémoire \texttt{MAP\_SHARED}, avec des indices \texttt{std::atomic}. En C++20, écrivez-le. Pourquoi les indices doivent-ils être sans verrou (lock-free), et comment le vérifier à la compilation~? Si le producteur est tué à une instruction quelconque, que peut voir le consommateur, et que ne peut-il jamais voir~? Comment un mutex dans la même mémoire se compare-t-il~? ThreadSanitizer ne suit pas deux processus, alors comment tester les ordres mémoire~?
\end{question}
